\documentclass[12pt]{article}
\usepackage[spanish]{babel}
\usepackage[utf8]{inputenc}
\usepackage[T1]{fontenc}
\usepackage{geometry}
\geometry{margin=2.5cm}

\title{Informe de preparacion de base de datos espacial}
\author{Proyecto LADM\_COL SINIC WebGIS BDE}
\date{\today}

\begin{document}

\maketitle

\section{Objetivo}
Preparar una base de datos espacial PostgreSQL/PostGIS para gestionar datos
catastrales de forma academica, tomando como referencia el enfoque
LADM\_COL SINIC V1.0.

\section{Base de datos}
La base se denomina \texttt{bde\_ladm\_sinic}. Se usan las extensiones
\texttt{postgis} y \texttt{pgcrypto}. No se usa \texttt{postgis\_topology}.

\section{Schemas}
La organizacion logica se realiza con los schemas \texttt{ladm\_sinic},
\texttt{catalogos} y \texttt{auditoria}. Las tablas principales del proyecto se
ubican en \texttt{ladm\_sinic}.

\section{Tablas principales}
Se implementan las tablas \texttt{interesado},
\texttt{unidad\_administrativa}, \texttt{derecho\_interesado},
\texttt{unidad\_espacial}, \texttt{topografia\_representacion} y
\texttt{cartografia\_catastral}.

\section{Componente espacial}
Las tablas \texttt{unidad\_espacial}, \texttt{topografia\_representacion} y
\texttt{cartografia\_catastral} incluyen campos de geometria PostGIS con SRID
EPSG:9377. Se crean indices GIST para mejorar consultas espaciales.

\section{Carga de datos}
El script de datos de prueba inserta registros suficientes para validar los
CRUD del backend y la visualizacion tabular del frontend.

\end{document}
